%! TEX root = master.tex
\chapter{3}{Rational and Irrational Numbers}

The rationals are closed under arithmetic, and a little parity is enough to
prove that $\sqrt2$ is not among them. Once one irrational number is known,
many more follow, including the golden ratio.

\section{Arithmetic of the rationals}\label{sec:rational-arithmetic}

%! TEX root = ../master.tex
\begin{theorem}[Closure of the rationals]\label{thm:rational-closure}
  If $p, q \in \Q$, then $p + q$, $p - q$ and $pq$ are rational, and so is
  $p / q$ when $q \neq 0$.
\end{theorem}

\begin{proof}
  Write $p = a/b$ and $q = c/d$ with $a, b, c, d \in \Z$ and $b, d \neq 0$.
  Then
  \[
    p \pm q = \frac{ad \pm bc}{bd}, \qquad
    pq = \frac{ac}{bd}, \qquad
    \frac pq = \frac{ad}{bc} \quad (c \neq 0),
  \]
  and in each case the numerator and denominator are integers, with a nonzero
  denominator.
\end{proof}


\section{The irrationality of \texorpdfstring{$\sqrt2$}{sqrt 2}}\label{sec:sqrt2}

%! TEX root = ../master.tex
\begin{theorem}[$\sqrt2$ is irrational]\label{thm:sqrt2-irrational}
  There are no integers $a$ and $b \neq 0$ with $(a/b)^2 = 2$.
\end{theorem}

\begin{proof}
  Suppose there are, and argue by contradiction. Replacing $a$ and $b$ by
  $-a$ and $-b$ if necessary, we may take $b > 0$. By well-ordering
  (\cref{ax:well-ordering}), choose such a pair with $b$ as small as
  possible. Then $a^2 = 2b^2$ is even, so $a$ is even by
  \cref{thm:square-parity}. Write $a = 2c$. Then $4c^2 = 2b^2$, so
  $b^2 = 2c^2$ is even, and $b = 2d$ is even too. Now $(c/d)^2 = 2$ with
  $0 < d < b$, which contradicts the choice of $b$.
\end{proof}


\begin{remark}
  The proof is a contradiction (\cref{sec:methods}). The well-ordering axiom
  lets us start from the smallest possible denominator. The parity argument
  then produces a smaller one. Without well-ordering, the same argument would
  only show that any such pair can be halved forever.
\end{remark}

%! TEX root = ../master.tex
\begin{theorem}[Rational and irrational]\label{thm:rational-plus-irrational}
  Let $q \in \Q$ and let $x \in \R \setminus \Q$. Then $q + x$ is irrational,
  and so is $qx$ if $q \neq 0$.
\end{theorem}

\begin{proof}
  If $q + x = p$ were rational, then $x = p - q$ would be rational by
  \cref{thm:rational-closure}. If $qx = p$ were rational with $q \neq 0$,
  then $x = p / q$ would be rational. Both contradict $x \notin \Q$.
\end{proof}


\begin{exercise}\label{ex:three-plus-root2}
  Show that $3 + \sqrt2$ and $\frac{\sqrt2}{2}$ are irrational.
\end{exercise}

\begin{solution}
  $3$ and $\frac12$ are nonzero rationals and $\sqrt2$ is irrational
  (\cref{thm:sqrt2-irrational}). By \cref{thm:rational-plus-irrational}, so
  are $3 + \sqrt2$ and $\frac12\sqrt2$.
\end{solution}

\section{The golden ratio}\label{sec:golden-ratio}

\begin{exercise}\label{ex:sqrt5}
  Prove that $\sqrt5$ is irrational, and hence that the golden ratio
  $\varphi = \frac{1 + \sqrt5}{2}$ is irrational.
\end{exercise}

\begin{proof}
  \emph{Step 1: if $5 \mid a^2$ then $5 \mid a$.} Divide $a$ by $5$, so that
  $a = 5q + r$ with $r \in \set{0, 1, 2, 3, 4}$. Then
  $a^2 = 5(5q^2 + 2qr) + r^2$, and $r^2 \in \set{0, 1, 4, 9, 16}$ leaves
  remainder $0, 1, 4, 4, 1$ on division by $5$. So $5 \mid a^2$ only when
  $r = 0$.

  \emph{Step 2: $\sqrt5$ is irrational.} Copy the proof of
  \cref{thm:sqrt2-irrational}. Suppose $a^2 = 5b^2$ with $b > 0$ as small as
  possible. Step 1 gives $a = 5c$. Then $b^2 = 5c^2$, so $b = 5d$ with
  $0 < d < b$ and $(c/d)^2 = 5$, a contradiction.

  \emph{Step 3: $\varphi$ is irrational.} $\frac{\sqrt5}{2}$ is irrational,
  and so is $\varphi = \frac12 + \frac{\sqrt5}{2}$, by two applications of
  \cref{thm:rational-plus-irrational}.
\end{proof}

\begin{remark}
  An irrational number may still be the root of a polynomial with integer
  coefficients. $\sqrt2$ is a root of $x^2 - 2$, and $\varphi$ is a root of
  $x^2 - x - 1$. Such numbers are \emph{algebraic}. Numbers that are not
  algebraic are \emph{transcendental}; $e$ (Hermite, 1873) and $\pi$
  (Lindemann, 1882) are the classical examples. So $\varphi$ is irrational
  but \emph{not} transcendental. Those two theorems are much deeper than
  anything in this chapter.
\end{remark}
